\subsection{Пружинки}
%%%%%%%%%%%%%%%%%%%%%%%%%%%%%%%%%%%%%%%%%%%%%%%%%%%%
%%%%%%%%%%%%%%%%%%%%%%%%%%%%%%%%%%%%%%%%%%%%%%%%%%%%

 \z[{\begin{wrapfigure}[5]{r}{0.2\textwidth}%
		\vspace{-15pt}%
		\begin{flushright}%
 			\includegraphics[width=0.18\textwidth]{4.3.1}%
		\end{flushright}%
	\end{wrapfigure}%
	}]%
	{Груз массы $M$ и шарик массы $m$ висят на трёх невесомых пружинах одинаковой жёсткости (см рис.). Верхняя пружина отрывается от шарика в точке $A$. Определите ускорение $\vec{a}$ (модуль и направление) шарика в начальный момент после отрыва.}
 %
{$a=\frac{M+m}{m}g\sqrt{2}$, вниз-влево}
%%%%%%%%%%%%%%%%%%%%%%%%%%%%%%%%%%%%%%%%%%%%%%%%%%%% 
 \z[{\begin{wrapfigure}[3]{r}{0.2\textwidth}%
		\vspace{-15pt}%
		\begin{flushright}%
 			\includegraphics[width=0.18\textwidth]{4.3.2}%
		\end{flushright}%
	\end{wrapfigure}%
	}]%
	{Система из трех одинаковых шаров, связанных одинаковыми пружинами, подвешена на нити. Нить пережигают. Найдите ускорения шаров сразу после пережигания нити.}
 %
{$a_\text{верх}=3g, a_\text{ниж} = 0$}
%%%%%%%%%%%%%%%%%%%%%%%%%%%%%%%%%%%%%%%%%%%%%%%%%%%% 
 \z[{\begin{wrapfigure}[5]{r}{0.2\textwidth}%
		\vspace{+20pt}%
		\begin{flushright}%
 			\includegraphics[width=0.2\textwidth]{4.3.3}%
		\end{flushright}%
	\end{wrapfigure}%
	}]%
	{Тело массы $m$ соединено двумя пружинами жесткости $k_1$ и $k_2$ с неподвижными стенками, пружины первоначально не деформированы. При возникших колебаниях наибольшее ускорение тела равно $a$. Найдите максимальное отклонение тела от положения равновесия и максимальные силы, с которыми пружины действуют на стенки.}
%
{$x_{\max}=\frac{ma}{k_1+k_2};$~$F_{1\max}=k_1x_{\max};$~$F_{2\max}=k_2x_{\max}$}
%%%%%%%%%%%%%%%%%%%%%%%%%%%%%%%%%%%%%%%%%%%%%%%%%%%%
 \z{Витую пружину с начальной длиной $l$, жёсткостью $k$ и массой $m$ свернули в кольцо и соединили концы. После этого её раскрутили с угловой скоростью $\omega$ вокруг оси, проходящей через центр кольца перпендикулярно его плоскости. Найдите радиус кольца $R$ как функцию от $\omega$. Диаметр витков пружины много меньше её длины.}
%
{$R=\frac{l}{2\pi-\frac{m\omega^2}{2\pi k}}$}
%%%%%%%%%%%%%%%%%%%%%%%%%%%%%%%%%%%%%%%%%%%%%%%%%%%% 
 \z[{\begin{wrapfigure}[6]{r}{0.3\textwidth}%
		\vspace{-15pt}%
		\begin{flushright}%
 			\includegraphics[width=0.3\textwidth]{4.3.5}%
		\end{flushright}%
	\end{wrapfigure}%
	}]%
	{ Найдите общий коэффициент жёсткости системы пружин, изображённой на рисунке, если внешняя сила прикладывается к верхней платформе в вертикальном направлении. Лестница, на которую опираются пружины, бесконечна. Все платформы при сжатии пружин сохраняют горизонтальное положение и не касаются ступенек лестницы. Каждая из платформ, кроме самой верхней, опирается на две пружины. Коэффициенты жёсткости всех пружин одинаковы и равны $k$, оси всех пружин вертикальны. Массой пружин и платформ можно пренебречь. }
%
{$k_\Sigma = k\sqrt{5}$}
%%%%%%%%%%%%%%%%%%%%%%%%%%%%%%%%%%%%%%%%%%%%%%%%%%%%